\section{Core tactics}\label{sec:05-tactics:02-core-tactics:1}

\texttt{fun hpq hp => hpq hp} proves \(P\to Q\to Q\) (given a function \(P\to Q\) and a proof of \(P\), apply the one to the other) in one shot, because the whole term was already obvious before writing it. Most goals are not like that: the term to write depends on the goal currently open, which itself changes as pieces of it get filled in. Writing a proof incrementally, goal by goal, is exactly what tactics do; term-mode is the destination, tactic-mode is how to get there when the destination isn't visible yet.

\begin{lstlisting}
theorem modus_ponens {P Q : Prop} : (P (*$\rightarrow$*) Q) (*$\rightarrow$*) P (*$\rightarrow$*) Q := by
  intro hpq hp
  exact hpq hp
\end{lstlisting}

The goal opens as \texttt{(P → Q) → P → Q}, a function type, matching exactly the shape \texttt{fun hpq hp => ...} above would need to produce. \href{https://lean-lang.org/doc/reference/latest/Tactic-Proofs/Tactic-Reference/}{\texttt{intro}} peels off one \(\lambda\)-binder at a time: \texttt{intro hpq} turns the goal \texttt{(P → Q) → P → Q} into the smaller goal \texttt{P → Q} with \texttt{hpq : P → Q} now a hypothesis in context, and \texttt{intro hp} repeats this once more. What is left, the goal \texttt{Q} with \texttt{hpq : P → Q} and \texttt{hp : P} in scope, is exactly what a finished term needs to supply. \href{https://lean-lang.org/doc/reference/latest/Tactic-Proofs/Tactic-Reference/}{\texttt{exact}} closes it by naming that term directly: \texttt{hpq hp} has type \texttt{Q}, matching the goal exactly, so the proof is done. Two tactics reconstruct, piece by piece, the same term \texttt{fun hpq hp => hpq hp} written in one line above.

\texttt{intro}/\texttt{exact} alone force writing every intermediate term by hand, even when a hypothesis already available is \emph{itself} a function that gets to the goal.

\begin{lstlisting}
theorem apply_example {P Q : Prop} (hpq : P (*$\rightarrow$*) Q) (hp : P) : Q := by
  apply hpq
  exact hp
\end{lstlisting}

The goal is \texttt{Q}, and \texttt{hpq : P → Q} is already in context. Writing \texttt{exact hpq hp} closes it in one step, but when the argument itself isn't obvious yet, \href{https://lean-lang.org/doc/reference/latest/Tactic-Proofs/Tactic-Reference/}{\texttt{apply}} works backward from the goal instead: \texttt{apply hpq} matches the goal \texttt{Q} against \texttt{hpq}'s conclusion, leaving a new goal for exactly the piece still missing, \texttt{P}, \texttt{hpq}'s premise. \texttt{exact hp} closes that. For a hypothesis with several premises, \texttt{f : A₁ → ⋯ → Aₙ → G}, \texttt{apply f} against goal \texttt{G} leaves \texttt{n} new goals, one per premise, the working-mathematician move ``by \texttt{f}, it remains to check the hypotheses of \texttt{f}.''

None of \texttt{intro}/\texttt{exact}/\texttt{apply} touch a hypothesis that is already known to equal something else, only the goal.

\begin{lstlisting}
theorem rw_example (a b : Nat) (h : a = b) : a + 1 = b + 1 := by
  rw [h]
\end{lstlisting}

The goal \texttt{a + 1 = b + 1} and the hypothesis \texttt{h : a = b} differ by exactly one substitution. \href{https://lean-lang.org/doc/reference/latest/Tactic-Proofs/Tactic-Reference/}{\texttt{rw}} performs it: \texttt{rw [h]} replaces every occurrence of \texttt{h}'s left-hand side with its right-hand side in the goal, turning \texttt{a + 1 = b + 1} into \texttt{b + 1 = b + 1}, closed automatically by \texttt{rfl} once both sides are syntactically identical.

The substitution \texttt{rw} performs is not limited to the goal.

\begin{lstlisting}
theorem rw_at_example (a b c : Nat) (h1 : a = b) (h2 : a + c = 10) : b + c = 10 := by
  rw [h1] at h2
  -- h2 is now : b + c = 10, exactly the goal
  exact h2
\end{lstlisting}

The goal here, \texttt{b + c = 10}, already matches a hypothesis, \texttt{h2}, up to the same substitution \texttt{h1} performs on the goal in the previous example. \texttt{rw [h1] at h2} runs the identical left-to-right rewrite, but targets the named hypothesis \texttt{h2} instead of the goal: every occurrence of \texttt{a} inside \texttt{h2} becomes \texttt{b}, turning \texttt{h2 : a + c = 10} into \texttt{h2 : b + c = 10}, which \texttt{exact h2} then closes directly. The direction and substitution rule are unchanged from ordinary \texttt{rw}; only the target moves from goal to hypothesis. The ring proofs of Chapter 10 use \texttt{rw [...] at h1}/\texttt{at h2} repeatedly for exactly this reason, to reshape a hypothesis into the form a later \texttt{exact} needs.

\begin{mathreading}

Each tactic corresponds to a standard proof move. \texttt{intro} discharges an implication/universal by the deduction theorem, to prove \(A \to B\), \emph{assume} \(A\) as a new hypothesis and prove \(B\). This is the \(\lambda\)-abstraction rule. \texttt{exact e} supplies a finished term, ``this is precisely our claim.'' \texttt{apply f} is backward chaining. \texttt{rw [h]} with \(h : a = b\) is substitution of equals for equals (Leibniz); every occurrence of \(a\) in the goal is replaced by \(b\), justified because \(a = b\) makes the old and new goals equivalent.

\end{mathreading}

\begin{quote}
Read more. ``Deduction theorem'' and ``\(\lambda\)-abstraction rule'' are two names for the same rule. ``Deduction theorem'' names the \(\Rightarrow\)-intro rule from natural deduction, stated in \hyperref[sec:04-propositions-and-proofs:02-logic-recap:1]{Chapter 4, Section 2}. ``\(\lambda\)-abstraction rule'' names the Curry--Howard reading of that same rule as a Lean \texttt{fun}, part of the typed \(\lambda\)-calculus formalized in \hyperref[sec:06-rigor-check:03-typing-rules-and-safety:1]{Chapter 6, Section 3}.
\end{quote}
